\documentclass[leqno, a4paper, 12pt]{jsarticle}
\usepackage{amssymb}
\usepackage{amsthm}
\usepackage{amsmath}
\usepackage{comment}
\usepackage{url}

\usepackage{enumitem}
%\usepackage{showkeys}


%定理環境 thmはあまり使わずpropをなるべく使う。
%主張は基本的に証明の中で使い、そのため番号を付けない。
%注意環境を付けてみた。
\newtheorem{thm}{定理}
\newtheorem{prop}[thm]{命題}
\newtheorem{cor}[thm]{系}
\newtheorem{lem}[thm]{補題}
\newtheorem{df}[thm]{定義}
\newtheorem{exam}[thm]{例}
\newtheorem{fact}[thm]{事実}
\newtheorem*{claim}{主張}
\newtheorem*{cau}{注意}
\newtheorem*{rmk}{注意}
\renewcommand{\proofname}{証明}
%矢印たち。
%\toは\rightarrowと同じ。\getsは\leftarrowと同じ。同値は\iff
\newcommand{\To}{\Longrightarrow}
\newcommand{\Gets}{\LongLeftarrow}
%黒板太字
\newcommand{\nn}{\mathbb{N}}
\newcommand{\zz}{\mathbb{Z}}
\newcommand{\qq}{\mathbb{Q}}
\newcommand{\rr}{\mathbb{R}}
\newcommand{\cc}{\mathbb{C}}
%無限大
\newcommand{\mgn}{\infty}
%差集合は\setminusで統一する。等号の否定は\neq
%真部分集合は\subsetneqq　偏微分は\partial
%ディスプレイスタイル。\dfracでディスプレイ分数
\newcommand{\dpst}{\displaystyle}

\title{論文メモ
\large{\\--- ボレル階層とか ---}
}
\author{ナンブ　キトラ}
\date{\today}
\begin{document}
\maketitle

本棚を整理したいので，
溜まっている論文のメモを書いて未来への手紙とすることにする．

以下の論文は
ボレル階層とかを調べようとして放置されていた論文の束だと思う．



論文
\cite{MR0138088}の用語をおさらいしよう．
一般に距離空間$X$が
\emph{絶対$\alpha$次ボレル階層}を持つ($\alpha$は可算順序数)
とは任意の距離空間$Y$
に$X$を位相的に埋め込んだ時に
$Y$の中で$X$が$\alpha$次のボレル集合になっている
時にいう．埋め込む先は距離空間というのが大事である．
この用語を踏まえた上で，よく知られた事実として
$X$が絶対$G_{\delta}$であることと
$X$がポーランドであることが同値である．
こういうのはブルバキに書いている．
ブルバキは一般性を志向している感じがするが
急にポーランド空間とか具体的な話を扱い始めるのでなんか物珍しいと思う．
論文
\cite{MR0138088}
では距離空間が絶対$F_{\sigma}$
になるための必要十分条件を導出している．
ちなみに$X$が可分の場合には絶対
$F_{\sigma}$と
$\sigma$コンパクト性が同値ということは知られていたようだ．
他にも大体以下のようなことを証明している．
\begin{thm}[\cite{MR0138088}]
以下が成り立つ．
\begin{enumerate}[label=\textup{(\arabic*)}]
\item 
距離空間$X$が絶対$F\cap G$空間(つまり開集合と閉集合の共通部分で絶対書かれる)
ための必要十分条件は$X$が局所コンパクトであることである．

\item 
$X$が絶対$F_{\sigma}$であるための
必要十分条件はそれが可算個の局所コンパクト(距離)空間の
和集合で表されること($\sigma$-locally compact)である．
\end{enumerate}
\end{thm}


論文
\cite{MR0219024}
は距離空間ではなく，
一般の位相空間の絶対閉性を扱っている．
位相空間$X$が絶対閉というのは
ハウスドルフ空間に埋め込んだ時に
必ず閉集合になるという意味で
これはコンパクト性の一般化と見做せる．
論文ではハウスドルフ空間は絶対閉空間に
閉集合として必ず位相的に埋め込めるとか
絶対閉包がproductiveになる必要十分条件とか
絶対閉空間の圏の射影対象や入射対象の研究をしているようだ．

論文
\cite{MR0260958}
ではタイトル通りボレル集合の和について研究している．
ここで和と言っているが
\textbf{和集合ではなく}，
$\rr$の二つの部分集合$A$と
$B$に対して$\rr$の足し算を使って定義される
\[
A+B=\{\, a+b\mid a\in A, b\in B\, \}
\]
である．
ざっくりいうと$A$と$B$がボレル集合だとしても
$A+B$がボレルとは限らなし，
例え$A$と$B$がそれぞれコンパクト，$G_{\delta}$
だとしても$A+B$がボレルとは限らないという例を構成している．コンパクトとか$G_{\delta}$というのは人間にとってかなり都合のいいボレル階層であることに注意しよう．
また，コンパクト集合と閉集合の和は閉集合になり，
閉集合は距離空間の中だと$G_{\delta}$集合であることにも注意しよう．この事実からコンパクト集合と$F_{\sigma}$集合の
和が$F_{\sigma}$になることがわかるので，
コンパクト集合と$G_{\delta}$の和というのはすでに知られていることと不明だったことのちょうど境目である．
エルデシュとストーンはフォンノイマンが構成した
$\rr$の連続体を持つ部分集合であって
$\qq$上代数的独立な具体的(ZFで構成できる)
なものを使って，
(まあこれはカントール集合に同相だったりするのだが)，
和がボレルにならないものを構成している．
また$\rr$以外の位相群でもできるっぽいことを注意している．


この種の話題はこれらの論文に限らない．
これらの論文が引用している論文や，
これらの論文を引用している論文を
確かめられたい．

%READ!
%myplain is the same to amsplain style. 
\nocite{*}
\bibliographystyle{myplaindoidoi}
\bibliography{bobo5.bib}



\end{document}